%HOMEWORK template for M 325K
\documentclass{article}
\headheight=7pt         \topmargin=14pt
\textheight=574pt       \textwidth=445pt
\oddsidemargin=18pt     \evensidemargin=18pt

%Begin package import

\usepackage{amsmath, amssymb, amsthm}
\usepackage{mathtools}
\usepackage{pdfpages}
\usepackage{tikz-cd}

%End package import
%Begin custom commands

\newcommand{\C}{\mathbb{C}}
\newcommand{\R}{\mathbb{R}}
\newcommand{\Q}{\mathbb{Q}}
\newcommand{\Z}{\mathbb{Z}}
\newcommand{\N}{\mathbb{N}}
\newcommand{\F}{\mathbb{F}}
\newcommand{\abs}[1]{\left\lvert #1\right\rvert}
\newcommand{\RM}[1]{\mathrm{#1}}
\newcommand{\BF}[1]{\textbf{#1}}
\newcommand{\e}{\epsilon}
\newcommand{\ceil}[1]{\lceil{#1}\rceil}
\newcommand{\floor}[1]{\lfloor{#1}\rfloor}
\newcommand{\rarr}[0]{\rightarrow}
\newcommand{\IM}[1]{\text{Im}(#1)}
\newcommand{\Hom}[0]{\text{Hom}}
\newcommand{\Pow}[1]{\text{Pow}(#1)}
\newcommand{\angles}[1]{\langle #1 \rangle}

%End custom commands
%Begin custom theorems

\newtheorem{innercustomgeneric}{\customgenericname}
\providecommand{\customgenericname}{}
\newcommand{\newcustomtheorem}[2]{%
\newenvironment{#1}[1]
{%
\renewcommand\customgenericname{#2}%
\renewcommand\theinnercustomgeneric{##1}%
\innercustomgeneric%
}
{\endinnercustomgeneric}
}

\newcustomtheorem{THRM}{Theorem}
\newcustomtheorem{LEM}{Lemma}
\newcustomtheorem{EX}{Exercise}
\newcustomtheorem{COR}{Corollary}
\newcustomtheorem{PROB}{Problem}
\newcustomtheorem{claim}{Claim}

\newenvironment{solution}
{\renewcommand\qedsymbol{$\blacksquare$}\begin{proof}[Solution]}
{\end{proof}}

\newenvironment{definition}
{\renewcommand\qedsymbol{$\blacksquare$}\begin{proof}[Definition]}
{\end{proof}}

%End custom theorems
\begin{document}

\title{Finite Abelian II}
\author{Arham Lodha}%put your name here
\date{March 19, 2024}%enter the due date here
\maketitle

\begin{PROB}{7}\label{Problem 7.}
\end{PROB}
\begin{proof}
    Let $R$ be a PID. $\forall a, b \in \R$, by definition of a PID, $(a, b) = (g)$ where $g = xa + yb$ where $x, y \in \R$. Similarly by the definition of a PID, $\exists h \in R \ni (x, y) = h$ thus $\exists u, v \in \R \ni x = hu$ and $y = hv$. Thus $g = xa + yb = hua + hvb = h(ua + vb)$. $k:= ua + vb$ and $k \in (a, b) = (g)$, thus $g | k$. However $g = hk \implies k | g$. Thus $g = k \implies h = 1$. Thus $(x, y) = (1)$. Thus $\exists p, q \in R, \ni xq + yp = 1$. Thus $\begin{bmatrix}
        x & -p \\ y & q
    \end{bmatrix}$ is invertible with determinant 1. Moreover, $\begin{bmatrix}
        a & b \\ * & *
    \end{bmatrix} \begin{bmatrix}
        x & -p \\ y & q 
    \end{bmatrix} = \begin{bmatrix}
        xa + yb & * \\ * & *
    \end{bmatrix} = \begin{bmatrix}
        g & * \\ * & *
    \end{bmatrix}$
    
    
    Suppose $(a)$ is maximal among the set of ideals generated by any element of $RAC$, for any possible invertible $R$ and $C$. Suppose $a$ is in the (1, 1) entry of $RAC$. 

    For the $i$-th  entry of the first row. Since we are in a PID, there is a Division Algorithm, so $a_{1, i} = qa_{1, 1} + r$. Apply the column operation, $c_{i}-qc_{1} \to c_{i}$. Thus the $0 < a_{1, i} = r < a_{1, 1}$, but since $a$ is the maximal ideal, $r = 0$. Thus $a_{1, 1} | a_{1, i}$. 
    
    For the $i$-th  entry of the first column. Since we are in a PID, there is a Division Algorithm, so $a_{i, 1} = qa_{1, 1} + r$. Apply the row operation, $r_{i}-qr_{1} \to r_{i}$. Thus the $0 < a_{i, 1} = r < a_{1, 1}$, but since $a$ is the maximal ideal, $r = 0$. Thus $a_{1, 1} | a_{i, 1}$. 

    Now apply the arguement in 4. 
\end{proof}

\bigskip

\begin{PROB}{8}\label{Problem 8.}
\end{PROB}
\begin{proof}
    Define the function $deg: \F[X] \to \Z^+$ which $\forall k \in \Z^+$, $\forall f_1, \dots, f_k \in \F$, takes $\sum_{i = 0}^k f_i x^i \to k$. $\forall g, h \in \F[X]$, $g = \sum_{i = 0}^{k_1}a_i x^i$ and $h = \sum_{i = 0}^{k_2} b_i x^i$. Thus $gh = \left(\sum_{i = 0}^{k_1} a_i x^i\right)\left(\sum_{i = 0}^{k_2} b_i x^i\right) = \sum_{i = 0}^{k_1k_2} c_i x^i$. Thus $deg(g) = k_1, deg(h) = k_2, deg(gh) = k_1k_2$. Thus $deg(h) \leq deg(gh) \geq deg(gh)$.
    
    $\forall n, d \in F[X] \ni d \neq 0$, suppose $deg(n) < deg(d)$ take $q = 0$ and $r = n$.
    
    Otherwise $\deg(n) \geq \deg(d)$, 
    
    let $n(x) = \sum_{k = 0}^{\deg(n)} a_k x^k$ and $d(x) = \sum_{k = 0}^{\deg(d)} b_k x^k$. 
    
    We can subtract from $n$, a suitible multiple of $d$ such that it eliminates the highest degree term of $n$. 
    
    So $r(x) = n(x) - \frac{a_{\deg(n)} x^{d - n}}{b_{\deg(d)}} m(x)$. $\deg(r) < \deg(f)$. If $\deg(r) \geq \deg(d)$ continue the process. 
    
    Eventually, $r(x) = n(x) - d(x) * \left(\frac{a_{\deg(n)}}{b_{\deg(d)}} x^{n - d} + \dots \right)$
    
    Thus $\deg(r) < \deg(d)$ and $q(x) = \left(\frac{a_{\deg(n)}}{b_{\deg(d)}} x^{n - d} + \dots \right)$. 

    Thus $F{X}$ has a division Algorithm and $\deg$ is a euclidean function.  
    
    Thus $F{X}$ is a euclidean domain.  


    $F \subset F[X]$ naturally. Let $M$ be a S-module. M is a abelian group which F acts on. Thus $M$ is a F-vector space. However $f(x) \in F[X]$ also acts on $M$. So one must define a function from $\phi: S \to Hom_{F}(M, M)$ which maps $f(x)$ to some linear map. 
    
    $\forall f(x) \in F[X], f(x) = \sum_{k = 0}^{n}f_{k}x^k$, $\forall k, f_{k} \to f_{k}I$, or scalar multiplication, because $F$ acts by scalar multiplication on any $F$ vector space. Thus $\phi(f(x)) = \sum_{k = 0}^{n}f_{k}I \circ \phi(x^k)$ 
    
    Moreover $X^k$ is the repeated multiplication of $X$ in $F[X]$ and must map to the repeated composition of $\phi(X)$ in $Hom_F(M, M)$ (ie $\phi(X) \circ ... \circ \phi(X)$ ) since $\phi$ is a ring homorphism. Thus $\phi(X^k)$ is completely determined by where $\phi(X)$ goes. Thus $\phi(X) $ uniquely determines $\phi$.  

    
    
    % $\alpha: Hom_{r}(S, Hom_{F}(M, M)) \to Hom_{F}(M, M)$ where $f \to f(X)$. $\forall g, h, a, b \in F[X]$, $\alpha(gh + ab) = (gh + ab)(X) = g(X)h(X) + a(X)b(X) = \alpha(g)\alpha(h) + \alpha(a)\alpha(b)$. Thus $\alpha$ is a ring homorphism. Let $\beta: Hom_{F}(M, M) \to Hom_{r}(S, Hom_{F}(M, M))$ where $A \to [f \to f(A)]$. $\forall A \in Hom_{F}(M, M)$, $(\alpha \circ \beta)(A) = \alpha([f \to f(A)]) = A$ and $\forall \gamma \in Hom_{r}(S, Hom_{F}(M, M)) $, $(\beta \circ \alpha)(\gamma) = \beta(\gamma(X)) = \gamma$. 

    % Thus $Hom_{r}(S, Hom_{F}(M, M)) \cong Hom_{F}(M, M)$. 
    
    % Thus $\phi$ is determined, and isomorphic, to $\phi(X)$.
    
    
\end{proof}

\bigskip

\begin{PROB}{9}\label{Problem 9.}
\end{PROB}
\begin{proof}
    $\implies$: Suppose $M$ is a S-module and $\exists f \in F[X]/\left\{ 0 \right\} \ni \forall m \in M, fm = 0 \implies f(A) = 0$. By Classification of S modules, $M = \bigoplus_{k = 1}^n \left(\frac{F[X]}{\left(p_k\right)}\right) \oplus F[X]^r$. $\forall k, \left(\frac{F[X]}{\left(p_k\right)}\right) \cong span(1, x, \dots, x^{\deg(p_k) - 1})$, thus $\left(\frac{F[X]}{\left(p_k\right)}\right)$ is finite dimensional and  $\bigoplus_{k = 1}^n \left(\frac{F[X]}{\left(p_k\right)}\right)$ is finite dimensional. 
    
    
    Assume the contrary, $r \neq 0$. Then take $\forall q_1, \dots q_r \in F[X]/\left\{ 0 \right\}$, then $m = (0, \dots, 0, q_1(x), \dots, q_r(x)) \in M$ and $f(x)(0, \dots, 0, q_1(x), \dots, q_r(x)) = (0, \dots, 0, f(x)q_1(x), \dots, f(x)q_r(x)) = f(A)m = 0$ so $\forall i, f(x)q_i(x) = 0$. But since $F[X]$ is a integral domain, $f(x)q_i(x) \neq 0$ since $f(x) \neq 0$ and $q_i(x) \neq 0 $. 

    Thus $r = 0$ by contradiction. Thus $M = \bigoplus_{k = 1}^n \left(\frac{F[X]}{\left(p_k\right)}\right)$ and thus is finite dimensional.
    
    
    $\impliedby:$ Suppose M is a finite dimensional $F$ vector space and a S-module. Thus $M \cong \bigoplus_{k = 1}^n \left(\frac{F[X]}{\left(p_k\right)}\right)$. Take $f = \prod_{k = 1}^{n}p_k$. 
    
    $\forall q \in F[X]$, $fq \in (p_i)$ for all $i$. Thus $\forall m \in M$, $fm = 0$, because each element of $m$ will be in the ideal of $(p_i)$ and thus will be 0.        
\end{proof}

\bigskip

\begin{PROB}{10}\label{Problem 10.}
\end{PROB}
\begin{proof}
    $\implies$: Suppose $F^n = M_{A} \cong M_{B} = F^n$ as $S$ modules. Thus there $\exists$ a $S$-linear isomorphism $\alpha: M_A \to M_B$, such that $\forall m \in M_A$, $\forall f(X) \in F[X]$ $\phi(f(A) m) = f(B) \phi(m)$. Thus the following commutative diagram is observed, 
    
    
    % https://q.uiver.app/#q=WzAsNCxbMCwwLCJGXm4iXSxbMiwwLCJGXm4iXSxbMCwyLCJGXm4iXSxbMiwyLCJGXm4iXSxbMCwxLCJcXGFscGhhIiwwLHsic3R5bGUiOnsidGFpbCI6eyJuYW1lIjoiaG9vayIsInNpZGUiOiJ0b3AifSwiaGVhZCI6eyJuYW1lIjoiZXBpIn19fV0sWzAsMiwiQSIsMl0sWzEsMywiQiJdLFsyLDMsIlxcYWxwaGEiLDIseyJzdHlsZSI6eyJ0YWlsIjp7Im5hbWUiOiJob29rIiwic2lkZSI6InRvcCJ9LCJoZWFkIjp7Im5hbWUiOiJlcGkifX19XV0=
\[\begin{tikzcd}
	{F^n} && {F^n} \\
	\\
	{F^n} && {F^n}
	\arrow["\alpha", hook, two heads, from=1-1, to=1-3]
	\arrow["A"', from=1-1, to=3-1]
	\arrow["B", from=1-3, to=3-3]
	\arrow["\alpha"', hook, two heads, from=3-1, to=3-3]
\end{tikzcd}\]

    Thus $B = \alpha A \alpha^{-1}$, and $A$ is conjugate to B. 
    
    
    $\impliedby$: Suppose A is conjugate to B. So there there is a isomorphism $\alpha: F^n \to F^n \ni B = \alpha A \alpha^{-1}$. Thus $\alpha \circ A = B \circ A$. Thus $\alpha \circ A^k = B^k \circ \alpha$. Thus $\alpha \circ f(A) = f(B) \circ \alpha$, $\forall f\in F[X]$. Thus $\phi$ is an isomorphism of $M_A$ and $M_B$ since $M_A \cong F^n$ and $M_B \cong F^n$.            
\end{proof}

\end{document}